\documentclass[11pt]{article}

\usepackage[margin=0.88in]{geometry}
\usepackage{amsmath,amssymb}
\usepackage{array}
\usepackage{microtype}
\usepackage{seqsplit}
\usepackage{longtable}
\usepackage{xcolor}
\usepackage{hyperref}
\usepackage{url}

\hypersetup{
  colorlinks=true,
  linkcolor=blue!55!black,
  citecolor=blue!55!black,
  urlcolor=blue!55!black
}
\setlength{\emergencystretch}{2em}
\newcommand{\method}{DreaMS-conditioned Hybrid MS2LDA}
\newcommand{\code}[1]{\texttt{#1}}
\newcommand{\hash}[1]{\texttt{\seqsplit{#1}}}
\newcolumntype{L}[1]{>{\raggedright\arraybackslash}p{#1}}

\title{\method:\\
Classical Topic Discovery with Semi-Amortized Spectral Inference}
\author{MS2LDA Development Team}
\date{6 August 2026}

\begin{document}
\maketitle

\begin{abstract}
MS2LDA discovers recurring fragment and neutral-loss patterns by treating
tandem mass spectra as documents in latent Dirichlet allocation (LDA).
This note specifies a two-stage hybrid. First, sparse classical variational LDA
discovers free topics by full-corpus expected-count updates, assisted by a
bounded DreaMS-conditioned topic--word prior. Second, topic discovery is frozen
permanently and a DreaMS document encoder is trained with a semi-amortized local
LDA objective: the exact local evidence lower bound (ELBO) is differentiated
through two unrolled coordinate updates. The encoder is never trained or used
inside topic discovery. The implementation exposes a small experiment-facing
interface, a one-way finalization lifecycle, and weights-only inference
checkpoints. Focused tests and fresh-encoder frozen-topic experiments support a
narrow limited-budget inference claim. A full-data, $K=1000$, leakage-safe
seed-42 MSnLib experiment additionally provides genuine held-out chemical
evidence. Under the primary full-spectrum, compound-balanced SOS diagnostic,
Tomotopy and long-refinement HybridLDA have 598 and 627 evaluable topics and
mean annotation-containment SOS of 0.6232 and 0.6293, respectively. This
supports chemically comparable motifs under a stricter split than the
single-seed, fitted-document analysis in the Nature Communications MS2LDA 2.0
paper. It does not establish universal chemical superiority. MS2LDA is a
research project at this stage; the result informs the preferred research
method rather than a production-system migration.
\end{abstract}

\section{Introduction}

Latent Dirichlet allocation represents a document as a mixture of latent
topics and each topic as a probability distribution over words
\cite{blei2003}. MS2LDA transfers this construction to tandem mass
spectrometry: spectra are documents, discretized fragments and neutral losses
are words, and topics are Mass2Motifs that may represent recurring molecular
substructures \cite{vanderhooft2016}. One spectrum can contain several motifs,
and one motif can recur across otherwise different molecules.

There are two distinct tasks. In \emph{discovery}, the topic--word
distributions are learned from a corpus. In \emph{decomposition}, a new
spectrum is assigned mixture weights over already learned topics. A network
that predicts membership in fixed motifs can accelerate decomposition, but it
cannot discover new motifs. A fully neural topic model can learn both parts,
but changes the optimizer and can suffer component collapse or unstable topic
identification \cite{srivastava2017}.

This reference takes a middle route. Free topics are updated from variational
LDA expected counts, while a frozen public spectral encoder supplies continuous
features to an empirical-Bayes word prior and, after discovery, a
semi-amortized new-document initializer. The two neural uses are deliberately
separated: pooled peak features can bias the bounded word prior during
discovery, whereas global spectrum embeddings train only the final local
inference network. The scientific question is whether these features add
useful chemical structure without replacing LDA's auditable topic update. A
full-data seed-42 benchmark now provides evidence for fast local inference and
chemically comparable topic discovery on benchmark-unseen compounds. This
document reports the complete design, disclosed post-hoc correction, and all
available outcomes without selecting favourable arms. The one-seed scope
matches the published MSnLib Tomotopy/SOS precedent, while a five-seed extension
remains useful for robustness rather than as an asymmetric requirement imposed
only on HybridLDA.

The contributions are:
\begin{enumerate}
  \item a frozen DreaMS feature extractor that returns aligned spectrum-level
        and contextual peak representations;
  \item one hybrid LDA model with classical topic discovery followed by a
        required semi-amortized inference-finalization stage;
  \item a deliberately small adapter surface for fitting, inference, topic
        access, and safe inference-only checkpoints; and
  \item a reproducible comparison with Tomotopy and the frozen DreaMS
        representation, including a completed full-data seed-42 execution and
        a deferred five-seed robustness extension.
\end{enumerate}

\section{Background}

\subsection{Classical and neural topic models}

Mean-field variational inference for LDA alternates local document updates
with global expected-count updates \cite{blei2003,hoffman2013}. Tomotopy uses
a collapsed Gibbs sampler rather than mean-field variational inference, but
fits the same broad LDA factorization. Both methods are stochastic and
label-invariant: different seeds can find different local solutions even when
all topics are used.

Neural topic models such as ProdLDA replace repeated local optimization with
an amortized encoder \cite{srivastava2017}. This makes inference fast and
allows continuous covariates, but creates an amortization gap: one shared
mapping need not return the best variational posterior for every document
\cite{cremer2018}. Semi-amortized inference reduces this gap by using the
encoder output as the start of local optimization \cite{kim2018}. Our model
uses this idea while retaining a conjugate LDA topic posterior.

\subsection{Pretrained DreaMS representations}

DreaMS is a transformer pretrained by self-supervision on millions of
unannotated tandem mass spectra \cite{bushuiev2026}. It emits
1,024-dimensional contextual representations for spectral peaks and a
spectrum-level representation at its precursor token. Its training objective
and attention architecture encode peak context and pairwise mass differences,
information absent from an MS2LDA count vector. We use the public
\code{DreaMS\_embedding} checkpoint without fine-tuning the transformer.

\section{Method}

\subsection{LDA core}

Let $x_{dv}$ be the count of feature $v$ in spectrum $d$, with $K$ motifs and
$V$ retained features. Conditional on the empirically estimated
$\boldsymbol\eta$, the model is classical LDA:
\begin{align}
  \boldsymbol\beta_k &\sim
    \operatorname{Dirichlet}(\boldsymbol\eta_k), \\
  \boldsymbol\theta_d &\sim
    \operatorname{Dirichlet}(\boldsymbol\alpha), \\
  z_{dn} &\sim \operatorname{Categorical}(\boldsymbol\theta_d), \\
  w_{dn} &\sim \operatorname{Categorical}
    (\boldsymbol\beta_{z_{dn}}).
\end{align}
The topic--word variational posterior is
$q(\boldsymbol\beta_k)=\operatorname{Dirichlet}(\boldsymbol\lambda_k)$,
and the document posterior is
$q(\boldsymbol\theta_d)=\operatorname{Dirichlet}(\boldsymbol\gamma_d)$.
For each nonzero feature count,
\begin{align}
  \phi_{dvk} &\propto \exp\left[
    \psi(\gamma_{dk})-\psi\left(\sum_j\gamma_{dj}\right)
    +\psi(\lambda_{kv})-\psi\left(\sum_u\lambda_{ku}\right)
    \right],\\
  \gamma_{dk} &= \alpha_k+\sum_v x_{dv}\phi_{dvk},\\
  \lambda_{kv} &= \eta_{kv}+\sum_d x_{dv}\phi_{dvk}.
\end{align}
The last equation is a full-corpus expected-count update. The neural
components do not replace it. The configured $\boldsymbol\alpha$ initializes
discovery rather than remaining fixed. After each training E-step, a guarded
Newton M-step maximizes its asymmetric Dirichlet objective using only
$E_q[\log\boldsymbol\theta_d]$ from training documents.

The implementation names the sparse count matrix \code{x}, the local
parameters \code{gamma}, the responsibilities \code{phi}, and the global
parameters \code{lambda\_posterior}. The method \code{beta\_mean()} returns
the posterior mean of $\boldsymbol\beta$. These names intentionally follow the
equations.

\subsection{Frozen spectrum and peak features}

For spectrum $d$, the frozen DreaMS checkpoint produces a global representation
$\mathbf s_d\in\mathbb R^{1024}$ and a contextual representation
$\mathbf h_{dp}\in\mathbb R^{1024}$ for each retained peak $p$. The extractor
uses the official DreaMS preprocessor and its top 100 peaks.

Contextual word features are constructed from training spectra only. For a
fragment word at mass $m_v$, we pool the state of the nearest peak within
0.02 Da. For a neutral-loss word at loss $\ell_v$, we pool the nearest peak to
precursor mass minus $\ell_v$. Repeated occurrences contribute according to
their counts. This gives a mean contextual vector $\bar{\mathbf h}_v$.
Unmatched words receive a zero contextual vector but retain their mass and
word-type features.

\subsection{Chemically structured topic--word prior}

The contextual vector, numeric mass, and feature type are combined into a
trainable normalized word representation:
\begin{equation}
 \mathbf u_v = \operatorname{normalize}\left[
 I_v W_h\operatorname{LN}(\bar{\mathbf h}_v)+
 g_m(\widetilde m_v)+\mathbf e_{\operatorname{type}(v)}
 \right].
\end{equation}
Here $I_v\in\{0,1\}$ records whether contextual peak evidence was observed, so
an unmatched word contributes no contextual projection. The scalar input is
$\widetilde m_v=\log(1+\max(m_v,0))/\log(2001)$. Here $g_m$ is a small
multilayer perceptron, and
$\mathbf e_{\operatorname{type}(v)}$ distinguishes fragments, losses, and
other tokens. Each topic has a trainable normalized vector $\mathbf t_k$.
The prior for topic $k$ is
\begin{align}
 p_{kv} &=
 \operatorname{softmax}_v\left(
 \frac{\mathbf t_k^\top\mathbf u_v}{T}\right),\\
 \eta_{kv}^{(e)} &= \eta+
 r_e\,\rho\,\frac{N}{K}p_{kv},
\end{align}
where $N$ is the number of training tokens. The configuration defaults are
$T=0.5$, $\rho=0.05$, and $r_e=\min(e/15,1)$, giving a 15-epoch warm-up.
Every topic therefore receives the same structured prior mass.

With row-normalized topic vectors collected in $\widetilde T$, the prior
parameters minimize
\begin{equation}
 -\frac{1}{KV}\sum_k \mathbb E_q[\log
   \operatorname{Dir}(\boldsymbol\beta_k;\boldsymbol\eta_k^{(e)})]
 +10^{-3}\frac{\|\widetilde T\widetilde T^\top-I\|_F^2}{K^2}.
\end{equation}
With the default configuration, one Adam update is taken per global epoch
through epoch 20 and the prior is then fixed. Automatic discovery convergence
is declared only when the relative $L_1$ changes in both
$\boldsymbol\lambda$ and $\boldsymbol\alpha$ are below $10^{-3}$ for three
consecutive fixed-prior epochs; 100 epochs is the default maximum. These are
configurable operational rules, not a proof of ELBO convergence. A caller may
deliberately stop discovery and finalize after any completed epoch, accepting
the risk of locking a premature topic state.

\subsection{DreaMS-conditioned document encoder}

The encoder is anchored to the final free topics. Define
\begin{align}
 \ell_{kv} &=
 \psi(\lambda_{kv})-\psi\left(\sum_u\lambda_{ku}\right),\\
 r_{vk} &= \operatorname{softmax}_k(\ell_{kv}),\\
 e_{dk} &= \sum_v\frac{x_{dv}}{N_d}r_{vk},
 \qquad N_d=\sum_v x_{dv}.
\end{align}
For an empty in-vocabulary document, $\mathbf e_d$ is defined as the uniform
topic vector. The spectrum representation follows the exact projection
\begin{equation}
 \mathbf p_d=\operatorname{LN}_2\!\left(
   \operatorname{GELU}\!\left[W_s\operatorname{LN}_1(\mathbf s_d)+\mathbf b_s\right]
 \right),
\end{equation}
which has 128 output dimensions by default. It is concatenated with
$\mathbf e_d$ and passed through two 256-unit ReLU hidden layers. The output is
a topic-logit residual:
\begin{align}
 \mathbf a_d &= g_\omega\left[
 \mathbf e_d;\mathbf p_d\right],\\
 \widehat{\boldsymbol\pi}_d &=
 \operatorname{softmax}\left(
   \log\!\left[\max(\mathbf e_d,\epsilon)\right]+\mathbf a_d\right),\\
 \boldsymbol\gamma_d^{(0)} &=
 \boldsymbol\alpha+N_d\widehat{\boldsymbol\pi}_d.
\end{align}
The final residual layer starts at zero, so the untrained network begins from
topic evidence supplied by LDA rather than an unrelated coordinate system. The
construction also preserves
$\sum_k\gamma_{dk}^{(0)}=\sum_k\alpha_k+N_d$ exactly.

Training VB starts from the retained $\boldsymbol\gamma_d$ from the preceding
epoch, not from the encoder. Consequently, global spectrum embeddings affect
new-document inference but not discovered topics. Only pooled contextual peak
features affect discovery, through the structured word prior. This is an
explicit separation of classical discovery and neural local inference, not a
ProdLDA-style end-to-end topic model. The encoder and spectrum projector are
neither evaluated nor optimized during discovery.

During discovery, document posteriors are refined with up to 50 sparse local
coordinate updates and tolerance $10^{-4}$ before the global expected-count
update. For a new spectrum, the encoder returns an immediate estimate, which
may be followed by standard local LDA updates. Refinement changes only that
document's topic mixture; it does not change learned motifs.

\subsection{Required semi-amortized inference finalization}

After the final topic update, \code{finalize\_inference()} performs the required
second stage exactly once. The complete topic posterior
$\boldsymbol\lambda$ and all structured-prior parameters are treated as fixed,
and a fresh Adam optimizer updates only the document encoder and spectrum
projector for 12 epochs by default. The implementation snapshots and checks the
topic posterior and every trainable structured-prior parameter bit-for-bit
before declaring finalization successful. If optimization fails, it restores
the encoder, random-generator state, and module mode. Topic discovery cannot
resume after successful finalization, and new-document inference and
checkpoint saving are unavailable beforehand.

For a candidate document posterior $\boldsymbol\gamma_d$, write
\begin{equation}
 \bar\theta_{dk}=\psi(\gamma_{dk})-
 \psi\left(\sum_j\gamma_{dj}\right).
\end{equation}
With the categorical responsibilities set to their coordinate optimum, the
encoder-dependent local sparse LDA ELBO is compactly
\begin{equation}
 \mathcal L_d(\boldsymbol\gamma_d;\boldsymbol\lambda)
 =-\operatorname{KL}\!\left[
   \operatorname{Dir}(\boldsymbol\gamma_d)\,\Vert\,
   \operatorname{Dir}(\boldsymbol\alpha)\right]
 +\sum_v x_{dv}\operatorname{LSE}_k
   \left(\bar\theta_{dk}(\boldsymbol\gamma_d)+\ell_{kv}\right).
 \label{eq:local-elbo-compact}
\end{equation}
Written out, up to constants that do not depend on document inference, this is
\begin{align}
 \mathcal L_d(\boldsymbol\gamma_d;\boldsymbol\lambda)
 ={}& \log\Gamma(\alpha_0)-\sum_k\log\Gamma(\alpha_k)
      -\log\Gamma(\gamma_{d0})+\sum_k\log\Gamma(\gamma_{dk}) \notag\\
 &+\sum_k(\alpha_k-\gamma_{dk})\bar\theta_{dk}
 +\sum_v x_{dv}\log\sum_k
   \exp\left(\bar\theta_{dk}+\ell_{kv}\right),
 \label{eq:local-elbo}
\end{align}
where $\alpha_0=\sum_k\alpha_k$ and
$\gamma_{d0}=\sum_k\gamma_{dk}$.  The final term is exactly the usual
expected assignment term plus categorical entropy, evaluated with a stable
log-sum-exp.

Starting from the encoder output $\boldsymbol\gamma_d^{(0)}$, two fixed
coordinate steps are unrolled:
\begin{align}
 \phi_{dvk}^{(t)}
   &=\operatorname{softmax}_k\left(
     \bar\theta_{dk}^{(t)}+\ell_{kv}\right),\\
 \gamma_{dk}^{(t+1)}
   &=\alpha_k+\sum_v x_{dv}\phi_{dvk}^{(t)}.
\end{align}
Equivalently, these equations define a differentiable map
$F_d(\boldsymbol\gamma)$ with
$\boldsymbol\gamma_d^{(1)}=F_d(\boldsymbol\gamma_d^{(0)})$ and
$\boldsymbol\gamma_d^{(2)}=F_d(\boldsymbol\gamma_d^{(1)})$. Back-propagation
passes through both maps and into $\boldsymbol\gamma_d^{(0)}$; it does not pass
into $\boldsymbol\lambda$ or the structured prior. The single supported encoder
objective is
\begin{equation}
 \mathcal J_{\mathrm{SA}}
 =-\frac{\sum_d\left[
       \mathcal L_d(\boldsymbol\gamma_d^{(2)};\boldsymbol\lambda)
       +0.1\,\mathcal L_d(\boldsymbol\gamma_d^{(0)};\boldsymbol\lambda)
     \right]}{\sum_d N_d}.
 \label{eq:semi-amortized-objective}
\end{equation}
For a uniformly shuffled minibatch of $B$ documents from a corpus of $D$
documents, its ELBO sum is multiplied by
$D/(B\sum_d N_d)$. This is the correctly scaled uniform-document minibatch
estimator of Equation~\ref{eq:semi-amortized-objective}. The implementation
uses random reshuffling, Adam, and gradient-norm clipping; later minibatches in
an epoch therefore see updated parameters and should not be described as
independent draws. Normalizing separately by each random minibatch's token
total would instead reweight short and long batches.
Epoch progress values are online minibatch diagnostics because the encoder
changes between updates, not a fresh evaluation of the final epoch parameters.
The small zero-step term retains useful encoder-only predictions. Training uses
exactly two updates with no adaptive stopping inside the differentiable graph,
making the objective and its back-propagation path fixed for every document.

At inference time, the default is the same two-update budget, although a caller
may request another exact number of local updates. Alternatively, a positive
tolerance makes that number a maximum batch-level budget and stops when
$\max_{d,k}|\gamma_{dk}^{(t+1)}-\gamma_{dk}^{(t)}|/
\max(|\gamma_{dk}^{(t)}|,\epsilon)$ falls below the threshold. Adaptive
stopping is never used in the differentiable training graph. Reported
training-document topic mixtures are also recomputed from
retained classical-VB state under the final topics, rather than from the
neural initializer.

\subsection{Training and inference sequence}

The implementation follows one explicit sequence:
\begin{enumerate}
  \item build a training-only vocabulary and sparse document--word matrix,
        then initialize the topic posterior from a seeded random simplex;
  \item refine each training document's $\boldsymbol\gamma_d$ and
        $\boldsymbol\phi_{dv}$ using the current topics;
  \item accumulate full-corpus expected topic--word counts, update the bounded
        structured prior during its declared window, and set
        $\boldsymbol\lambda$ to prior plus expected counts;
  \item continue until the caller stops or stationary-prior convergence is
        declared;
  \item permanently freeze the final topics and train the encoder through two
        differentiable local updates using
        Equation~\ref{eq:semi-amortized-objective}; and
  \item for an unseen spectrum, compute the encoder initialization and
        optionally apply local VB without changing $\boldsymbol\lambda$.
\end{enumerate}
No dense document-by-vocabulary tensor is constructed. Each batch contains
only padded nonzero word identifiers and counts.

\subsection{Implementation and compatibility boundary}

Conceptually the system is one fitted \code{HybridLDAModel} plus one frozen
external feature extractor. \path{ms2lda_hybrid/_variational.py} contains the
sparse local LDA equations and ELBO, \path{ms2lda_hybrid/_core.py} contains the
neural modules and global topic update, and \path{ms2lda_hybrid/model.py} owns
the lifecycle and checkpoint boundary. \path{ms2lda_hybrid/dreams_features.py}
contains the lazy DreaMS loader, aligned HDF5 cache, and train-only peak-to-word
pooling. Pooling requires unique document identifiers and rejects any order
mismatch with cached DreaMS rows. There is no second model implementation.

\code{HybridLDAModel} provides only the experiment-facing operations needed to
add documents, run classical discovery, finalize inference, inspect topics,
infer new-document mixtures, and save or load inference checkpoints. Some names
mirror the existing model API where their semantics agree, but this is not a
drop-in Tomotopy replacement: in particular, it does not expose sampled
token-topic assignments. It is not wired into the existing CLI, dashboard, or
visualizer. This is an intentional research-interface boundary, not evidence
against using HybridLDA as the preferred method in subsequent research.

The version-3 reference checkpoint can be written only after successful
finalization. It contains the topic posterior, neural parameters,
configuration, vocabulary, and explicit finalized state. It deliberately
excludes training documents, local variational state, optimizer state, and
random-number-generator state and is therefore inference-only. It is loaded
with PyTorch's weights-only mode under PyTorch 2.6 or newer. DreaMS's legacy
Lightning checkpoints are accepted only after their exact public SHA-256
hashes and static pickle globals have been verified; the restricted loader
then allowlists only the required DreaMS architecture and metadata types.
Feature provenance is stored in the DreaMS cache rather than the model
checkpoint, so callers must keep compatible artifacts together.

A distinct version-1 training-state format is available only for controlled
experiment resume. It remains weights-only but additionally contains the
free-topic posterior, retained gamma matrix, discovery or encoder optimizer,
random states, convergence counters, and histories. It omits training spectra
and feature arrays, so restore first reconstructs those immutable inputs and
then verifies their frozen context hash, configuration, vocabulary, document
count, sparse-matrix shape, nonzero count, and token total. This format never
changes the version-3 inference artifact or permits discovery after encoder
training has begun.

\section{Validation status}

The focused tests verify configuration invariants, sparse count handling, the
local LDA updates and ELBO, monotonic coordinate refinement, differentiable
two-step refinement, conservation of token and Dirichlet mass, asymmetric
alpha estimation, fixed topic-word prior mass, joint alpha/topic convergence,
deterministic same-seed discovery and finalization, absence of encoder updates
during discovery, bitwise frozen-topic, alpha, and topic-word-prior invariance
during finalization, the one-way lifecycle, adaptive
inference, strict weights-only checkpoint round trips, feature caching,
identifier-aligned peak-to-word pooling, immutable vocabulary accessors, and
the shared spectral-word parser. Additional tests interrupt and resume both
training stages and require tensor-exact agreement with uninterrupted controls;
they also test context rejection, two-generation retention, and fallback from
a deliberately corrupt latest checkpoint.

\subsection{Frozen-topic inference experiments}

The fixed objective was first explored in a deterministic oracle-topic
simulator and then rerun on seeds 23, 29, and 31 from fresh seeded encoders. A
benchmark-only posterior-mean regression comparator received the same topics,
initial encoder, documents, minibatch order, and 12 training epochs. The
common per-document reference was the maximum ELBO among 100-step solves from
the symmetric initializer and both trained encoder initializers, followed by
one diagnostic continuation step. The
semi-amortized objective reduced the two-step ELBO gap by 13.37\%, 20.04\%, and
14.15\%, respectively. Their arithmetic mean was 15.86\%; the relative
reduction between the across-seed mean gaps was 15.84\%. Across-seed mean
encoder-only observed-token NLL rose by 0.086\%, while two-step NLL was slightly
lower. This simulator
isolates inference behavior; its oracle topics and latent-derived embeddings
are not chemical evidence.

The synthetic driver is
\path{benchmarks/semi_amortized_inference.py}. The regression comparator lives
only in \path{benchmarks/inference_baselines.py}; it is not exposed by the
reference package or serializable as a finalized checkpoint. The driver can
write all per-seed results as JSON. This experiment supports a narrow claim:
from a fresh seeded encoder and fixed topics, the final semi-amortized encoder
is a better limited-budget document-inference initializer than this comparator.
It does not show that discovered motifs are more chemically meaningful, more
stable, or superior to another topic-discovery method.

\section{Full leakage-safe MSnLib validation}
\label{sec:msnlib-validation}

\subsection{Question, source data, and frozen scale}

The benchmark question is whether semi-amortized \method{} can infer a topic
mixture close to a long local-variational solution, at close to encoder-only
speed, for spectra from genuinely unseen chemical structures, while discovering
motifs with chemical coherence comparable to ordinary Tomotopy LDA.  The
benchmark is implemented as the single staged driver in
\path{benchmarks/msnlib_validation/}. This repository is a research project;
the comparison selects and characterizes research methods rather than testing a
production migration.

The source files are the positive-mode MSnLib case study and deposited results
from Torres Ortega et al.\ \cite{torres2026}, together with the positive-mode
Spec2Vec model, embedding matrix, and SQLite reference library named by that
paper. They were downloaded from Zenodo record 20179680
(doi:10.5281/zenodo.20179680). Their immutable paths, byte sizes, and SHA-256
hashes are recorded in \path{configs/full-msnlib-k1000.json}; the external
asset inventory additionally records the archive hashes, Zenodo checksums,
extraction roots, and construction audit. The released MGF contains 41,568
source spectra. Applying the published preprocessing parameters retains 38,888
spectra; 2,680 spectra fail the minimum of five retained peaks. This is the
full paper-compatible eligible set, not a down-sampled benchmark. Every profile
uses $K=1000$ topics. The completed laptop execution uses seed 42, matching the
single seed in the deposited Nature-paper notebook; the unexecuted robustness
extension would add seeds 43--46. The original five-seed confirmatory
configuration is retained as an audit artifact, but the known seed-42 test
result means it can no longer create pristine confirmatory evidence.

The preprocessing interval is 0--2000 Da, with at most 1,000 peaks, relative
intensity 0.01--1.0, two decimal places, and \code{min\_df=3}. Each retained
peak creates an intensity-weighted fragment word and, when positive, a neutral
loss word. Both word types derived from one physical peak remain one indivisible
group during document completion. The ordinary LDA initial document-topic
prior is $\alpha=0.6$, and the topic-word prior is $\eta=0.1$. Both methods
optimize their document-topic prior from training spectra only. Tomotopy uses
50-iteration checkpoints, at most 5,000 iterations as in the deposited MSnLib
notebook, and the frozen ten-checkpoint 0.005 relative-perplexity stopping rule.
The report rejects a seed that reaches this maximum without satisfying the
stopping rule. Standard held-out Tomotopy inference uses 100 iterations. The
five-seed profile requests one deterministic worker; the completed seed-42
laptop profile requests six workers in partition mode and is explicitly not
bitwise reproducible, like the paper's automatic multicore run.

\subsection{Executed seed-42 parameter snapshot}

Table~\ref{tab:executed-parameters} is the complete reader-facing parameter
snapshot for the corrected result. The authoritative machine-readable values
remain the committed seed-42 configuration and the frozen protocol lock; no
parameter below was selected from chemical labels.

\begin{table}[ht]
\centering
\small
\begin{tabular}{L{0.29\linewidth}L{0.63\linewidth}}
Component & Executed value \\
\hline
Data and words & positive mode; 0--2000 Da; 5--1000 retained peaks; relative
intensity 0.01--1.0; two decimal places; fragment plus positive neutral-loss
word per peak; \code{min\_df=3}, \code{min\_cf=0}, \code{rm\_top=0} \\
Scale & 41,568 input spectra; 38,888 eligible spectra; 21,233 training-only
words; $K=1000$ \\
Split and completion & scaffold-group 70/10/20 split, seed 42; connectivity
fallback for empty scaffolds; observed peak-group fraction 0.5, completion
seed 42 \\
Common priors & $\alpha_k=0.6$ initialization; $\eta=0.1$; asymmetric
$\boldsymbol\alpha$ re-estimated from training documents \\
Tomotopy discovery & seed 42; 6 workers; partition parallelism; maximum 5,000
iterations; 50-iteration steps; convergence window 10; relative-perplexity
tolerance 0.005; stopped after 750 iterations \\
Tomotopy inference & 100 local iterations; one inference thread \\
Hybrid discovery & seed 42; batch size 32; 4 PyTorch intra-operation threads;
maximum 250 epochs; tolerance $10^{-3}$ for both topic and alpha relative
change; patience 5; stopped after 149 discovery epochs \\
Hybrid finalization & 12 encoder epochs; two unrolled local-VB updates in the
training objective; two retained rotating checkpoints \\
Held-out Hybrid arms & encoder only (0 updates), encoder plus 2 updates, and
long refinement of 50 updates; one inference thread \\
Topic diagnostics & top 10 words per topic; corpus-active threshold 0.001;
document-active threshold 0.01; training-document co-occurrence NPMI \\
MAG and SOS & top 20 motif words; search $k=1000$; 10 unique molecules;
Spec2Vec clustering cosine 0.9; MACCS consensus 0.8; dominant-topic,
compound-balanced primary SOS; probability $\geq0.5$ sensitivity \\
Latency & 128 frozen test spectra; five repeats; cached-feature and end-to-end
scopes reported separately \\
\end{tabular}
\caption{Complete parameter snapshot for the executed seed-42 run. The
250-epoch ceiling was a safety limit; training stopped when both frozen global
criteria satisfied patience.}
\label{tab:executed-parameters}
\end{table}

\subsection{Primary structural split and leakage controls}

The primary split assigns complete Bemis--Murcko scaffold groups to 70\%
training, 10\% validation, and 20\% test partitions by a deterministic
descending-group-size algorithm with seed 42. Repeated spectra sharing a
connectivity InChIKey therefore cannot cross partitions. Acyclic structures,
for which a Murcko scaffold is empty, fall back to connectivity identity rather
than being pooled into one meaningless empty-scaffold group. The frozen split
contains 27,222 training, 3,889 validation, and 7,777 test spectra, with zero
cross-partition connectivity-key or declared split-group overlap. Non-empty
scaffold groups stay intact; the 228 acyclic spectra use the connectivity
fallback described above. A random spectrum split is not used as primary
evidence.

\begin{table}[ht]
\centering
\small
\begin{tabular}{lrrrr}
Partition & target fraction & spectra & connectivity identities & split groups \\
\hline
Training & 0.70 & 27,222 & 26,916 & 20,014 \\
Validation & 0.10 & 3,889 & 3,854 & 2,856 \\
Test & 0.20 & 7,777 & 7,695 & 5,702 \\
\end{tabular}
\caption{Frozen structural split. A split group is a complete non-empty
Bemis--Murcko scaffold or, for an empty scaffold, one connectivity identity.
Both split-group and connectivity overlaps across partitions are zero.}
\label{tab:frozen-split}
\end{table}

Only training spectra construct the 21,233-word vocabulary and pooled
contextual peak representations. Chemical names, SMILES, InChIKeys, and MAG
annotations are absent from the model-facing assignment file and are not used
in fitting, convergence decisions, or reference-budget selection. For every
validation and test document, half of the physical peak groups are selected by
a spectrum-ID-keyed deterministic hash. The DreaMS representation is recomputed
using only those observed groups, whose intensities and word multiplicities are
renormalized after the split; neither is copied from or scaled by the complete
held-out spectrum. Completion words outside the training vocabulary are counted
and reported as OOV rather than silently assigned probability.

MAG is run only after all topic models finish. Every validation and test
connectivity key is removed from the 4.37-GB SQLite/Spec2Vec search library
before the FAISS index is built. The retained index stores an explicit mapping
to original database rows. An exclusion audit covers every retained row, and
every returned search hit is checked again before deserialization. Thus neither
the evaluation compound nor another spectrum of the same compound can be a MAG
reference. The MAG stage refuses a data root different from the frozen lock and
checks the topic and test-mixture hashes before applying chemical identities.
This is stricter than the deposited full-corpus MSnLib analysis,
whose reference library may overlap its evaluated compounds.

\subsection{Compared methods and frozen inference selection}

Every method receives identical training documents, training vocabulary,
topic count, split, and seed. The four topic-model arms are:
\begin{enumerate}
  \item ordinary Tomotopy LDA with standard 100-iteration held-out inference;
  \item HybridLDA encoder-only inference, \code{infer(iter=0)};
  \item the same encoder followed by exactly two local-VB updates; and
  \item the same frozen topics with a validation-selected long local-VB budget.
\end{enumerate}
The long budget is selected without reading test outcomes. For each trained
seed, validation mixtures after 25 and 50 updates are compared. Fifty updates
are accepted only when the median cosine is at least 0.999 and the fifth
percentile at least 0.99; otherwise the frozen maximum of 100 is checked against
50. The run fails instead of calling the result ``near-converged'' if the
maximum misses either threshold. The chosen 50- or 100-step budget is then
applied unchanged to test spectra. Hybrid discovery permits at most 30 global
epochs in the original five-seed profile. The completed seed-42 continuation
raises only that safety ceiling to 250 and stops after 149 discovery epochs,
followed by 12 inference-encoder epochs. All convergence decisions use training
or validation quantities, never chemical labels or test likelihood.

A raw global-DreaMS nearest-neighbour analysis is included as a non-topic-model
baseline. It searches the training partition only and reports DreaMS cosine and
Morgan-fingerprint similarity for each test spectrum. It cannot supply topic
coherence or document likelihood and is labelled accordingly.

\subsection{Frozen metrics and machine-readable outputs}

Predictive quality uses document-completion negative log likelihood per
in-vocabulary token:
\begin{equation}
 \operatorname{NLL}=-\frac{1}{N_{\mathrm{held}}}
 \sum_{d,v}x^{\mathrm{held}}_{dv}\log\left(
 \sum_k\widehat\theta_{dk}\widehat\beta_{kv}\right).
\end{equation}
Lower NLL is better. For encoder-only and two-step inference, cosine and
Jensen--Shannon convergence and the NLL gap are reported against the long
solution. Topic diagnostics comprise document- and corpus-level active-topic
counts, top-ten word diversity, training-document word-co-occurrence NPMI, and
cross-seed stability after full-vocabulary cosine-optimal Hungarian matching.
The same matching reports same-seed Tomotopy--Hybrid cosine and top-word
Jaccard overlap. All matching quantities are similarities, not correctness
scores.

Timing has two scopes over the same frozen 128-spectrum subset and five repeats:
cached-feature model inference, and end-to-end inference. Hybrid end-to-end
timing includes DreaMS preprocessing and extraction; Tomotopy end-to-end timing
includes word-to-document construction because it does not use DreaMS. Training
runtime, whole-test inference runtime, feature-extraction runtime, MAG runtime,
and process peak RSS are saved separately.

Chemical evaluation applies the deposited MAG procedure (the top 20 motif
words, ten unique molecules, Spec2Vec cosine 0.9, and an 0.8 MACCS consensus
threshold) to both topic models. MAG optimization then retains only motif
features shared with the recommended reference spectra. A topic without any
optimized feature has no usable MAG annotation and is separated from a topic
that has an annotation but no associated held-out spectrum.

Likelihood and chemistry deliberately use different test representations.
Document completion infers from its frozen observed peak groups and scores the
remaining groups. Chemical coherence recomputes DreaMS and topic mixtures from
all physical peak groups in the held-out spectrum; chemical structures remain
withheld until inference has finished. The primary post-hoc SOS diagnostic
assigns each test spectrum to its highest-probability topic. This rank-based
rule uses every spectrum once and avoids an absolute probability cutoff; it is
not invariant to every calibration or prior difference. The deposited paper's
topic-probability threshold of 0.5 is retained unchanged as a labelled
sensitivity analysis, not tuned to the observed Hybrid output.
Within each topic, the primary SOS mean gives each held-out connectivity
identity one vote, preventing replicate spectra from being treated as
independent chemical evidence. A spectrum-weighted value is retained as an
explicit sensitivity result.

The deposited downstream analysis notebook computes SOS using RDKit Daylight
fingerprints, a 0.9 consensus threshold, and intersection bits divided by
annotation bits. The supplementary methods instead state that the denominator
is the smaller fingerprint, while the benchmark MAG configuration uses MACCS
and a 0.8 consensus threshold. The present comparison keeps that frozen
MACCS/0.8 setting identical between Tomotopy and Hybrid and reports both
denominator definitions. Published SOS values are therefore contextual, not a
numerically comparable baseline. Matched Tomotopy--Hybrid annotations receive a
MACCS-consensus Jaccard agreement score. These quantities measure coherence and
agreement; they do not establish which individual peaks are genuine fragments
or noise. In particular, MSnLib cannot prove that the HybridLDA background
component is experimental background.

The preferred independent chemical endpoint is recovery of manual
motif--spectrum annotations. No such independent annotation file is present in
the frozen Zenodo inputs, so the report records this endpoint as unavailable and
does not substitute the compound structures or deposited automated MAG labels
as ground truth. The deposited 1,000-topic result is matched only as
same-dataset contextual similarity. The paper reports 432 topics removed during
MAG optimization, 145 without associated molecules, and 423 remaining; its
listed MACCS bins (158 high, 176 intermediate, and 78 low) sum to 412 rather
than 423 \cite{torres2026}. This unreconciled difference is retained explicitly
rather than silently repaired.

\subsection{How this benchmark differs from the Nature Communications MS2LDA 2.0 paper}

The Nature Communications paper establishes the relevant published precedent,
but it did not ask the same generalization question. Its deposited
\path{Benchmark_MAG_MSn.ipynb} fits one $K=1000$ Tomotopy model with seed 42,
5,000 requested iterations, \code{parallel=3}, and automatic worker count.
\path{Analysis_MSnLib.ipynb} loads that single saved model and computes one SOS
analysis from its fitted documents with the 0.5 membership heuristic. Thus a
five-seed result is useful additional robustness evidence, but it is not a fair
prerequisite imposed only on HybridLDA. Table~\ref{tab:nature-comparison}
separates exact carryovers from deliberate changes.

\small
\begin{longtable}{L{0.18\linewidth}L{0.35\linewidth}L{0.37\linewidth}}
\caption{Relationship to the deposited Nature Communications MS2LDA 2.0
methods. ``This benchmark'' refers to the completed corrected seed-42
experiment.}\label{tab:nature-comparison}\\
Dimension & Nature-paper MSnLib analysis & This benchmark \\
\hline
\endfirsthead
Dimension & Nature-paper MSnLib analysis & This benchmark \\
\hline
\endhead
Scientific aim & Demonstrate large-scale MS2LDA 2.0 motif discovery and MAG
annotation. & Compare Tomotopy and HybridLDA topic discovery and unseen-spectrum
inference under identical inputs. \\
Corpus and $K$ & Full positive MSnLib corpus after the released preprocessing;
$K=1000$. & Same released input, preprocessing, and $K=1000$; 38,888 eligible
spectra, without down-sampling. \\
Seed and fit & One Tomotopy fit, seed 42, 5,000 requested iterations,
\code{parallel=3}, automatic workers. & One completed fit per method, seed 42.
Tomotopy uses partition parallelism with 6 workers and stops at its frozen
criterion; Hybrid uses 4 intra-operation threads. Neither fit is claimed
bitwise reproducible. \\
Evaluation population & SOS molecules are obtained from the same documents used
to fit the saved LDA model; there is no structural train/test partition. &
Training, validation, and test are disjoint scaffold groups. SOS uses 7,777 test
spectra from 7,695 connectivity identities unseen during benchmark fitting. \\
Fitted state & Vocabulary, topics, and document mixtures use the full fitted
corpus. & Vocabulary, pooled DreaMS features, transformations, topics, and
learned priors use training spectra only. Chemical labels remain unavailable to
fitting and model selection. \\
Spectrum--topic association & Every fitted document--topic pair with probability
$\geq0.5$ is associated; repeat compounds may contribute repeatedly. & Primary
SOS assigns each full held-out spectrum to its strongest topic and gives each
connectivity identity one vote within a topic. The unchanged 0.5 rule and
spectrum-weighted scoring are sensitivities. \\
MAG reference library & Full deposited Spec2Vec reference library; no held-out
evaluation partition exists to exclude. & All validation/test connectivity
identities are removed before indexing: 53,426 of 426,280 rows excluded and
zero held-out rows retained. \\
MAG eligibility & Published analysis uses MAG optimization before the reported
chemical analysis. & The correction likewise requires a nonempty
MAG-optimized motif; a clustered topic alone is not called annotated. \\
Fingerprint/SOS arithmetic & Benchmark configuration uses MACCS consensus 0.8;
the downstream notebook recomputes RDKit Daylight consensus 0.9. The notebook
and supplement disagree on the denominator. & MACCS/0.8 is frozen identically
for Tomotopy and Hybrid. Both denominator definitions are reported, so the
within-benchmark comparison is fair but not an exact numerical reproduction of
the paper's downstream SOS values. \\
Inference question & No encoder-only, two-step, long-refinement, document
completion, or DreaMS nearest-neighbour comparison. & Adds those arms and
baselines without using them as chemical ground truth. \\
\end{longtable}
\normalsize

The principal methodological strengthening is therefore the structural
held-out design and MAG exclusion audit, not a larger seed count. On the SOS
criterion used to compare topic methods, the present one-seed result meets the
same replication depth as the publication while evaluating the harder question
of benchmark-unseen compounds. Its small Hybrid--Tomotopy difference still
supports comparability rather than universal superiority.

\subsection{Protocol freeze, execution, and output schema}

Before any test metric is inspected, the driver hashes all inputs and benchmark
source, validates the full MGF, freezes split and completion assignments,
records Git branch/commit/dirty state and dependency versions, and writes a
self-hashed \code{protocol.lock.json}. Subsequent stages refuse changed inputs,
manifests, or source. Runs are resumable by seed and method and execute each
model in a fresh process, so peak memory is meaningful. Generated embeddings,
models, the filtered MAG index, and result artifacts live outside Git. The
committed \path{output-schema.json} defines every small and large artifact, and
\code{execution\_log.jsonl} records exact commands, executables, Conda
environment names, timestamps, and completion states.

The canonical environments remain separate. \code{ms2lda-hybrid} supplies
patched PyTorch, DreaMS, HybridLDA, and Tomotopy; the five-seed Tomotopy profile
is deterministic, whereas the executed seed-42 profile explicitly is not. The
historical run used the local \code{MS2LDA\_v2} environment for the released
Spec2Vec/MAG stack. Its executed dependency versions are now captured portably
in \path{environment-msnlib-mag.yml} under the environment name
\code{ms2lda-msnlib-mag}. Separation avoids the incompatible NumPy, Python,
PyTorch, MatchMS, and RDKit constraints of DreaMS and legacy MAG. The complete
fresh-acquisition and environment-construction commands are:
{\footnotesize
\begin{verbatim}
DATA=/path/to/ms2lda-msnlib-validation/zenodo/20179680

python scripts/download_msnlib_validation_assets.py --data-root "$DATA"
conda env create -f environment-hybrid.yml
conda env create -f environment-msnlib-mag.yml
conda run -n ms2lda-msnlib-mag \
  python -m pip install --no-deps -e .
\end{verbatim}
}

The historical \path{full-msnlib-k1000.json} remains frozen for audit and must
not now be executed or reported as confirmatory. A future five-seed run remains
valuable for run-to-run robustness, but it needs a new explicitly post-hoc
evidence-scope label, all five seeds without favourable-run selection,
\code{--test-results-inspected}, and a safety ceiling capable of reaching the
unchanged training-only convergence rule. No earlier feature, model, or result
artifact may be treated as output of that future protocol. Implementing this
new scope is deferred because it is not required to reproduce the completed
seed-42 result and silently relabelling the old confirmatory configuration would
damage its audit value.

\subsection{Full-data single-seed indicative execution}

The five-seed robustness experiment remains a planned extension, but its runtime
is not practical on the available laptop. The executed seed-42 profile therefore
retains the complete set of 38,888 eligible spectra, the same structural split
and training-only vocabulary, $K=1000$, seed 42, all Hybrid inference arms,
leakage-filtered MAG, and the raw-DreaMS baseline. It cannot estimate
run-to-run stability, but its one-seed depth matches the Nature-paper Tomotopy
and SOS analysis.

An earlier seed-42 Hybrid execution used the scalar value $\alpha=0.6$ as a
fixed value for each of 1,000 topics. Its total document-prior concentration
was therefore 600, whereas ordinary Tomotopy had optimized its document-topic
prior during training. This mismatch strongly regularized Hybrid test mixtures
towards diffuse topic usage and confounded the likelihood, convergence, and
chemical-coherence comparisons. All numerical Hybrid results and active run
artifacts from that execution have been removed; they are not evidence for or
against the method.

Hybrid now treats the configured $\alpha=0.6$ as an initialization and
re-estimates an asymmetric vector using only training-document variational
posteriors after every discovery E-step. For $D$ training documents and
$s_k=\sum_d E_q[\log\theta_{dk}]$, the optimized part of the variational
objective is
\begin{equation}
 {\cal L}(\boldsymbol\alpha)
 =D\left[\log\Gamma\!\left(\sum_k\alpha_k\right)
 -\sum_k\log\Gamma(\alpha_k)\right]
 +\sum_k(\alpha_k-1)s_k .
\end{equation}
A guarded Newton update preserves strictly positive finite values and accepts
only objective-nondecreasing steps. Discovery convergence requires both
$\boldsymbol\lambda$ and $\boldsymbol\alpha$ to satisfy the frozen relative
change threshold. The learned vector is then frozen with the topics before
encoder fitting and held-out inference.

The completed model run remains immutable. The corrected chemical endpoint is
an explicitly post-hoc child run so the earlier half-spectrum SOS output cannot
be mistaken for the corrected metric:
{\footnotesize
\begin{verbatim}
REPO=/path/to/MS2LDA
DATA=/path/to/ms2lda-msnlib-validation/zenodo/20179680
ROOT=/path/to/ms2lda-msnlib-validation/runs/\
indicative-msnlib-k1000-seed42
SOURCE="$ROOT/continuation"
RUN="$SOURCE/chemical-correction"
CFG=benchmarks/msnlib_validation/configs/\
indicative-msnlib-k1000-seed42-chemical-correction.json

cd "$REPO"

conda run -n ms2lda-hybrid python -m benchmarks.msnlib_validation \
  freeze-chemical-correction --config "$CFG" --data-root "$DATA" \
  --run "$RUN" --repo-root "$PWD" --source-run "$SOURCE" \
  --reason "correct full-spectrum SOS, MAG optimization, rank-based \
association, and compound weighting after diagnosis"
conda run -n ms2lda-hybrid python -m benchmarks.msnlib_validation \
  reuse-core-artifacts --source-run "$SOURCE" --run "$RUN"

export DATA RUN REPO
/bin/zsh -lc '
  set -e
  cd "$REPO"
  conda run --no-capture-output -n ms2lda-hybrid \
    python -m benchmarks.msnlib_validation run-core --run "$RUN"
  conda run --no-capture-output -n ms2lda-hybrid \
    python -m benchmarks.msnlib_validation run-chemical-inference --run "$RUN"
  conda run --no-capture-output -n ms2lda-hybrid \
    python -m benchmarks.msnlib_validation _run-raw-dreams --run "$RUN"
  conda run --no-capture-output -n MS2LDA_v2 \
    python -m benchmarks.msnlib_validation run-mag \
      --run "$RUN" --data-root "$DATA"
  conda run --no-capture-output -n ms2lda-hybrid \
    python -m benchmarks.msnlib_validation report --run "$RUN"
' >> "$RUN/unattended.log" 2>&1
\end{verbatim}
}
The indicative configuration uses six Tomotopy workers in partition mode, four
PyTorch intra-operation threads for Hybrid discovery and encoder fitting, and
one thread for all held-out inference and latency measurements. The initial
run reached its 50-epoch safety ceiling: the topic-word relative change was
$1.99\times10^{-4}$, below the frozen $10^{-3}$ tolerance, while the
training-only alpha relative change was 0.0141. The driver therefore retained
the checkpoint and stopped before encoder finalization, MAG, or reporting.
The disclosed continuation changes only the maximum safety ceiling, from 50 to
250 total epochs. It preserves the $10^{-3}$ tolerance, five-epoch patience,
model and optimizer states, variational factors, random-number-generator state,
and every data and evaluation setting. It resumes epoch 50 and still stops as
soon as both global criteria remain satisfied for five epochs.

\subsection{Results and discussion}

The corrected chemical diagnostic completed for all 38,888 filtered spectra,
including 7,777 structurally held-out test spectra representing 7,695
connectivity identities. The post-hoc correction reused the exact completed
models; it changed neither topics nor inference parameters. It recomputed the
test representation from all peak groups, applied MAG motif optimization, and
used compound-balanced scoring. All declared mixture, MAG, and report hashes
verified, and the filtered MAG library retained zero held-out reference rows.
The corrected protocol SHA-256 is
\hash{855abc25bcdef08a8bd0958f9985cd1c40c088009ad841b38dfd767eabf65435};
its unchanged core models originate from protocol
\hash{6c223207d7199d1467cb96fbe5d90f1b82cf88abd9a2a005433524747c4080a7}.
The compact committed checkpoint manifest preserves these identifiers and the
hashes of the uncommitted large result artifacts.

Table~\ref{tab:core-results} records the completed predictive and systems
diagnostics. They answer the semi-amortized inference question but are not
chemical endpoints. Two local updates increase the median held-out mixture
cosine from 0.9501 to 0.9996 relative to 50-step refinement, while end-to-end
latency changes only from 33.08 to 34.19 ms per spectrum. Thus two-step
inference is close to the long-refinement mixture at near-encoder speed on this
run. Encoder-only completion NLL is 9.7577, essentially Tomotopy's 9.7569;
two-step and long-refinement NLL are higher. This is not contradictory: local
VB optimizes the observed-half variational objective, not held-out-token NLL,
so the long mixture is an optimization reference rather than predictive ground
truth.

\begin{table}[ht]
\centering
\footnotesize
\begin{tabular}{lrrrrrr}
Inference arm & NLL/token & cosine & cached ms & end-to-end ms & test s & active topics \\
\hline
Tomotopy standard & 9.7569 & -- & 58.63 & 58.74 & 449.97 & 363 \\
Hybrid encoder & 9.7577 & 0.9501 & 1.88 & 33.08 & 1.95 & 361 \\
Hybrid encoder + 2 VB & 10.0766 & 0.9996 & 2.96 & 34.19 & 11.64 & 356 \\
Hybrid long VB (50) & 10.1915 & 1.0000 & 19.14 & 50.83 & 169.23 & 357 \\
\end{tabular}
\caption{Completed seed-42 document-completion and timing results. Lower NLL
is better. Cosine is the median per-spectrum similarity to the long Hybrid
mixture; 1.0000 for the long arm is by definition. Latencies are medians over
five repeats on 128 frozen test spectra. End-to-end Hybrid timing includes
DreaMS extraction. ``Test s'' is vectorized inference over all 7,777 test
spectra and is not directly interchangeable with the per-spectrum latency
protocol.}
\label{tab:core-results}
\end{table}

Hybrid training, including resumed discovery and encoder finalization, took
65,680.4 s (18.24 h) with 3.20 GiB peak RSS; Tomotopy took 2,426.5 s
(40.44 min) with 1.59 GiB. Hybrid's final document-prior concentration was
$\sum_k\alpha_k=9.1117$, versus 6.8065 for Tomotopy and 600 at initialization.
The topic models use similar numbers of held-out corpus-active topics, but
Hybrid has higher top-ten word diversity (0.8387 versus 0.8148) and worse
training co-occurrence NPMI ($-0.4744$ versus $-0.2997$). After optimal
same-seed topic matching, Tomotopy--Hybrid mean cosine is 0.4768 and mean
top-word Jaccard is 0.0816; these are similarities between label-invariant
solutions, not correctness scores. Cross-seed stability is unavailable with
one seed. The raw-DreaMS training-neighbour baseline has mean and median Morgan
similarity 0.2207 and 0.1707; its exact connectivity and scaffold match rates
are both zero, as required by the split. None of these secondary diagnostics
is substituted for the SOS chemical endpoint below.

\begin{table}[ht]
\centering
\small
\begin{tabular}{llrrrrr}
Method & inference & MAG & rank SOS & containment & smaller-FP & 0.5 SOS \\
\hline
Tomotopy & standard & 607 & 598 & 0.6232 & 0.6617 & 333 \\
HybridLDA & encoder & 632 & 602 & 0.6333 & 0.6665 & 191 \\
HybridLDA & encoder + 2 VB & 632 & 621 & 0.6298 & 0.6607 & 57 \\
HybridLDA & long VB (50) & 632 & 627 & 0.6293 & 0.6603 & 64 \\
\end{tabular}
\caption{Post-hoc seed-42 chemical-coherence diagnostic. MAG is the
number of topics retaining at least one optimized motif feature. Rank SOS is
the number of those topics receiving at least one held-out compound when each
spectrum is assigned to its highest-probability topic. Containment and
smaller-FP are compound-balanced mean SOS under the two declared denominator
definitions. The 0.5 SOS column is the SOS-evaluable topic count under the
unchanged historical probability cutoff and is a sensitivity result only.}
\label{tab:corrected-sos}
\end{table}

The corrected evaluation reverses the earlier impression that Hybrid produced
far fewer chemically usable topics. Before MAG eligibility is considered, 988
Hybrid long-refinement topics and 987 Tomotopy topics are the strongest topic
for at least one full test spectrum. After MAG optimization and
compound-balanced association, their SOS-evaluable counts are similarly 627
and 598. Hybrid also has slightly more MAG-annotatable topics (632 versus 607).
Thus the large old coverage gap was not a shortage of used Hybrid topics.

The mechanism is probability sharpness. The median largest probability in a
full test mixture is 0.215 for Hybrid long refinement and 0.310 for Tomotopy.
Only 90 Hybrid topics ever cross 0.5, compared with 519 Tomotopy topics; after
requiring both a MAG annotation and an associated compound, only 64 and 333
topics remain. The fixed cutoff therefore compares the numerical sharpness of
two inference algorithms as well as chemical coherence. It is retained for
historical transparency but is not a defensible primary cross-method coverage
metric.

On the rank-based diagnostic, mean annotation-containment SOS is 0.6232 for
Tomotopy and 0.6293 for Hybrid long refinement. Encoder-only Hybrid is 0.6333
and the two-step arm is 0.6298. The smaller-fingerprint definition gives the
same qualitative picture. Compound-balanced and spectrum-weighted values
differ by at most 0.0002, so repeated spectra do not drive this seed's result.
These small numerical differences show comparable chemical coherence; they do
not establish superiority. The evaluation has one non-reproducible multicore
seed, the rank-based endpoint was introduced after diagnosing the failed
cutoff, and the models discover different label-invariant topics.

These are genuine held-out MSnLib chemical-coherence measurements, distinct
from the software smoke tests. The correction is explicitly post-hoc and is not
a numerical reproduction of the paper's downstream RDKit/0.9 SOS analysis. No
independent manual motif--spectrum annotations are present in the frozen
Zenodo assets, so that stronger external endpoint remains unavailable. However,
the publication also reports one seed and one SOS analysis. Under that fair
evidence standard, the stricter held-out result supports HybridLDA as a credible
preferred research method with chemical coherence comparable to Tomotopy. It
does not establish that the small numerical advantage is universal.
\section{Limitations}

The hybrid retains LDA's non-convex, label-invariant topic discovery, so
different seeds can still reach different solutions. DreaMS is trained to
represent whole-spectrum molecular context, whereas a Mass2Motif is intended
to capture a smaller recurring pattern. A useful global representation does
not therefore guarantee a useful motif prior. The scaffold split establishes
that validation and test compounds are unseen by the benchmark-fitted topic
models, transformations, and HybridLDA document encoder. It does not establish
that they were absent from the much larger corpus used to pretrain the public
DreaMS checkpoint, because a complete spectrum-level pretraining manifest is
not available for an identity audit. Consequently, any claim of unseen-chemical
generalization is conditional on the fixed representation and must be stated as
``unseen during benchmark fitting,'' not ``provably unseen during all upstream
pretraining.''

The structured-prior mass, warm-up, and temperature are fixed design choices
that have not yet been validated. The two-step objective and 0.1 zero-step
weight were selected during development; the current synthetic experiment is
therefore a mechanism check, not independent chemical evidence.
The corrected real-data execution uses only seed 42 and
non-bitwise-reproducible multicore fitting, so it cannot quantify run-to-run
topic stability. The historical SOS endpoint's fixed 0.5 membership cutoff is
retained as a sensitivity result but is demonstrably not calibrated across
these inference algorithms. The rank-based, compound-balanced diagnostic
avoids that absolute cutoff but was introduced after inspecting the failed
endpoint and considers only each spectrum's strongest topic. A future
untouched-dataset study should prespecify it and preferably add a mass-weighted
secondary-motif analysis. Repeating seeds 43--46 on the existing split would
measure robustness but would not undo the disclosed post-hoc endpoint change.
Finalization intentionally locks topic discovery, so fitting on a premature
topic state cannot be repaired without training a new model. The model is not
an end-to-end neural topic model, and it intentionally exposes only a focused
research interface. Finally, a passing software test suite and frozen-topic
inference improvement cannot replace external chemical validation; the SOS
experiment, rather than production-readiness criteria, supplies the chemical
evidence considered here.

\section{Reproducibility}

The self-contained reference surface consists of
the \path{ms2lda_hybrid/} package, focused tests, the synthetic and MSnLib
benchmark modules, \path{environment-hybrid.yml}, and this document.
The package is included in built wheels with optional dependencies exposed by
the \code{hybrid} extra, while remaining isolated from the existing
\path{MS2LDA/} import and workflow boundary.
The selective use of the preserved evaluation branches is audited in
\path{benchmarks/msnlib_validation/PROVENANCE.md}: only the deterministic
Tomotopy construction precedent and historical MACCS SOS arithmetic were
reimplemented, with no archive merge, private metadata, data, or generated
result copied.

For review of the completed result, commit
\hash{88b47d51b1f045e95810aeafafa042834d63372b} is the exact executed-source
checkpoint: every benchmark and \path{ms2lda_hybrid/} path matches the SHA-256
recorded by the frozen real-run \code{code\_manifest.json}. The original
protocol lock records its earlier parent commit plus the complete dirty-tree
listing, because the source was frozen before this checkpoint commit existed.
The small committed
\path{docs/hybrid_lda_seed42_checkpoint.json} records the protocol lineage,
split counts, key SOS results, leakage audit, large-artifact hashes, and claim
boundary. Large models, DreaMS caches, FAISS indices, and full generated reports
remain outside Git as required; their hashes permit integrity checking when the
artifact directory is shared separately.

The public inputs can be reconstructed without local history using the
standard-library, resumable downloader. The historical result can then be
reviewed at its exact-source commit, or a scientifically equivalent fresh
seed-42 run can be started in a new output directory:
{\footnotesize
\begin{verbatim}
git clone git@github.com:joewandy/MS2LDA.git
cd MS2LDA
git switch --track origin/hybrid-lda-msn-validation

DATA=/path/to/ms2lda-msnlib-validation/zenodo/20179680
RUN=/path/to/ms2lda-msnlib-validation/runs/seed42-fresh
CFG=benchmarks/msnlib_validation/configs/\
indicative-msnlib-k1000-seed42-chemical-correction.json

python scripts/download_msnlib_validation_assets.py --data-root "$DATA"
conda env create -f environment-hybrid.yml
conda env create -f environment-msnlib-mag.yml
conda run -n ms2lda-msnlib-mag \
  python -m pip install --no-deps -e .

git checkout 88b47d51b1f045e95810aeafafa042834d63372b
conda run -n ms2lda-hybrid python -m benchmarks.msnlib_validation \
  validate-inputs --config "$CFG" --data-root "$DATA"
conda run -n ms2lda-hybrid python -m benchmarks.msnlib_validation \
  preflight --config "$CFG" --data-root "$DATA"
conda run -n ms2lda-hybrid python -m benchmarks.msnlib_validation \
  freeze --config "$CFG" --data-root "$DATA" --run "$RUN" \
  --repo-root "$PWD" --test-results-inspected
conda run -n ms2lda-hybrid python -m benchmarks.msnlib_validation \
  run-core --run "$RUN"
conda run -n ms2lda-hybrid python -m benchmarks.msnlib_validation \
  run-chemical-inference --run "$RUN"
conda run -n ms2lda-hybrid python -m benchmarks.msnlib_validation \
  _run-raw-dreams --run "$RUN"
conda run -n ms2lda-msnlib-mag \
  python -m benchmarks.msnlib_validation run-mag \
  --run "$RUN" --data-root "$DATA"
conda run -n ms2lda-hybrid python -m benchmarks.msnlib_validation \
  report --run "$RUN"
\end{verbatim}
}
The fresh run will not be bitwise identical to the historical Tomotopy model
because both use multicore partition training. It reconstructs the frozen
scientific comparison; exact historical numbers require the separately hashed
model artifacts. This limitation is inherited from the published multicore
Tomotopy setup and is disclosed rather than hidden.

The environment uses Python 3.11 and pins NumPy 1.25.0, patched PyTorch 2.6.0,
and
DreaMS at Git commit
\hash{dbec3a0b514a99e5056cfccde4559fda8cfe8129}. Its required architecture
package is pinned at
\hash{d35b34b0caee38348098cda852bff30d8b25cd25}. Installing DreaMS with
\code{--no-deps} is deliberate: DreaMS metadata pins vulnerable PyTorch 2.2.1
and a moving \code{msml@main} requirement. The YAML instead supplies the
embedding-path dependencies, patched PyTorch 2.6.0, and the exact compatible
architecture commit. The real public embedding path is tested with this
override.

From the repository root, the checkout-local reference can be tested with:
{\footnotesize
\begin{verbatim}
conda env create -f environment-hybrid.yml
conda activate ms2lda-hybrid
python -m pip install --no-deps \
  "git+https://github.com/pluskal-lab/DreaMS.git@dbec3a0b514a99e5056cfccde4559fda8cfe8129"
python -m pytest -q tests/test_hybrid_lda.py
python -m benchmarks.semi_amortized_inference
\end{verbatim}
}

Before importing either dependency, extractor construction verifies the
PEP~610 installed-package commits against both the pinned DreaMS and legacy
architecture commits. It then downloads missing public checkpoints without
loading them and checks their exact hashes before restricted deserialization.
The backbone and embedding head are loaded directly on the requested device;
the upstream convenience loader is not used because it otherwise selects the
default CUDA device whenever CUDA is available. DreaMS caches record both
dependencies' expected and installed commits and package versions, both public
checkpoint hashes, embedding dimension, and peak limit. Model loading does not
enforce a cache match. Finalized reference checkpoints contain no spectra and
cannot resume training; the separate experiment-only training state described
above can resume only when its external frozen inputs pass the context audit.

\section{Conclusion}

\method\ is a focused reference implementation of conjugate free-topic
discovery followed by semi-amortized local inference. Expected counts discover
topics, the bounded structured prior biases their word distributions, the
training-only empirical-Bayes update estimates the asymmetric document-topic
prior, and only after those quantities are frozen does the document encoder
learn through the local LDA ELBO. The corrected full-data seed-42 experiment is
complete. On structurally held-out compounds, long-refinement HybridLDA and
Tomotopy yield 627 versus 598 SOS-evaluable topics and mean
annotation-containment SOS of 0.6293 versus 0.6232. The publication likewise
uses one seed and one SOS analysis, while the present split and MAG exclusion
are stricter. The result therefore supports HybridLDA as a chemically
comparable and credible preferred method for further MS2LDA research. It does
not prove that the small numerical difference is universal, that individual
peaks are correct fragments, or that manual chemical validation is unnecessary.

{\small
\begin{thebibliography}{10}

\bibitem{blei2003}
D. M. Blei, A. Y. Ng, and M. I. Jordan.
Latent Dirichlet allocation.
\emph{Journal of Machine Learning Research}, 3:993--1022, 2003.

\bibitem{hoffman2013}
M. D. Hoffman, D. M. Blei, C. Wang, and J. Paisley.
Stochastic variational inference.
\emph{Journal of Machine Learning Research}, 14:1303--1347, 2013.

\bibitem{vanderhooft2016}
J. J. J. van der Hooft, J. Wandy, M. P. Barrett, K. E. V. Burgess, and
S. Rogers.
Topic modeling for untargeted substructure exploration in metabolomics.
\emph{Proceedings of the National Academy of Sciences},
113(48):13738--13743, 2016.
doi:10.1073/pnas.1608041113.

\bibitem{srivastava2017}
A. Srivastava and C. Sutton.
Autoencoding variational inference for topic models.
In \emph{Proceedings of the 5th International Conference on Learning
Representations}, 2017.

\bibitem{cremer2018}
C. Cremer, X. Li, and D. Duvenaud.
Inference suboptimality in variational autoencoders.
In \emph{Proceedings of the 35th International Conference on Machine
Learning}, PMLR 80:1078--1086, 2018.

\bibitem{kim2018}
Y. Kim, S. Wiseman, A. Miller, D. Sontag, and A. Rush.
Semi-amortized variational autoencoders.
In \emph{Proceedings of the 35th International Conference on Machine
Learning}, PMLR 80:2678--2687, 2018.

\bibitem{bushuiev2026}
R. Bushuiev, A. Bushuiev, R. Samusevich, C. Brungs, J. Sivic, et al.
Self-supervised learning of molecular representations from millions of tandem
mass spectra using DreaMS.
\emph{Nature Biotechnology}, 44:630--640, 2026.
doi:10.1038/s41587-025-02663-3.

\bibitem{torres2026}
L. R. Torres Ortega, J. Dietrich, J. Wandy, H. Mol, and
J. J. J. van der Hooft.
Large-scale discovery and annotation of substructure patterns in mass
spectrometry profiles.
\emph{Nature Communications}, 2026.
doi:10.1038/s41467-026-75038-0.

\end{thebibliography}
}

\end{document}
